% GENERATED FILE. Edit canonical lessons, then run npm run generate:books.
% canonical-source-hash: fnv1a64:198198e7128d4f7f
% canonical-lessons: UR-C413-kothi, UR-C413-jhonpri, UR-C413-ahata, UR-C413-chhappar, UR-C413-basana

\chapter{Houses Large and Small}
\label{ch:ur-houses-large-and-small}
\hlchaptermodality{413}

\begin{chapteropening}
\textbf{By the end of this chapter:} \emph{I can describe houses, from a hut with a thatch to a large house in its compound.}

Complete the last lesson of chapter 413: I can describe houses, from a hut with a thatch to a large house in its compound.
\end{chapteropening}

\section[\texorpdfstring{\ur{کوٹھی} (koṭhī)}{کوٹھی (koṭhī)}]{\ur{کوٹھی} (koṭhī) — a large house; a bungalow}

\label{lesson:UR-C413-kothi}



\begin{quote}
Before the new one: say the Urdu for the countryside, then the Urdu for next door.
\end{quote}

\subsection*{You'll want to know: \ur{کوٹھی}}
\textbf{\ur{کوٹھی}} — \emph{koṭhī} — \textquotedblleft{}a large house; a bungalow\textquotedblright{}.

\begin{quote}
\ur{امیر} \ur{لوگ} \ur{بڑی} \ur{کوٹھیوں} \ur{میں} \ur{رہتے} \ur{ہیں۔}
\end{quote}

\begin{quote}
\emph{amīr log baṛī koṭhiyoṅ meṅ rahte haiṅ.}
\end{quote}

\begin{quote}
Rich people live in big houses.
\end{quote}

Feminine.

\begin{cousinweb}
\textbf{\ur{کوٹھی}} (\emph{koṭhī}) is inherited: CDIAL lists the Hindi cousin \emph{koṭhī}, \textquotedblleft{}a granary, a large house\textquotedblright{}, under Sanskrit \emph{kóṣṭha}, \textquotedblleft{}a storeroom\textquotedblright{} (CDIAL 3546).
\end{cousinweb}

\begin{grammarlens}[title={What you've built}]
One more word for getting around.
\end{grammarlens}

\subsection*{Your turn}
\begin{itemize}
\raggedright
  \item \emph{Say it:} \emph{koṭhī}
  \item \emph{Say it:} \emph{koṭhī}, once more
  \item \emph{Say it:} \emph{amīr log baṛī koṭhiyoṅ meṅ rahte haiṅ.} — Rich people live in big houses.
  \item \emph{Recall:} say the Urdu for the countryside, then the Urdu for next door, then say \emph{koṭhī} again
\end{itemize}

\begin{tcolorbox}[breakable,colback=teal!4,colframe=teal!35!black,title={Before you move on}]
What does \ur{کوٹھی} mean? (\textquotedblleft{}A large house; a bungalow\textquotedblright{}.) Say it once more.
\end{tcolorbox}
\section[\texorpdfstring{\ur{جھونپڑی} (jhoṅpṛī)}{جھونپڑی (jhoṅpṛī)}]{\ur{جھونپڑی} (jhoṅpṛī) — a hut}

\label{lesson:UR-C413-jhonpri}



\begin{quote}
Before the new one: say the Urdu for next door, then the Urdu for a large house.
\end{quote}

\subsection*{You'll want to know: \ur{جھونپڑی}}
\textbf{\ur{جھونپڑی}} — \emph{jhoṅpṛī} — \textquotedblleft{}a hut\textquotedblright{}.

\begin{quote}
\ur{دریا} \ur{کے} \ur{کنارے} \ur{ایک} \ur{جھونپڑی} \ur{ہے۔}
\end{quote}

\begin{quote}
\emph{daryā ke kināre ek jhoṅpṛī hai.}
\end{quote}

\begin{quote}
There is a hut on the riverbank.
\end{quote}

Feminine.

\begin{cousinweb}
\textbf{\ur{جھونپڑی}} (\emph{jhoṅpṛī}) is inherited: CDIAL lists the Hindi cousin \emph{jhõpṛī}, \textquotedblleft{}a hut\textquotedblright{}, under \emph{jhuppa}, a form it reconstructs (CDIAL 5403).
\end{cousinweb}

\begin{grammarlens}[title={What you've built}]
One more word for getting around.
\end{grammarlens}

\subsection*{Your turn}
\begin{itemize}
\raggedright
  \item \emph{Say it:} \emph{jhoṅpṛī}
  \item \emph{Say it:} \emph{jhoṅpṛī}, once more
  \item \emph{Say it:} \emph{daryā ke kināre ek jhoṅpṛī hai.} — There is a hut on the riverbank.
  \item \emph{Recall:} say the Urdu for next door, then the Urdu for a large house, then say \emph{jhoṅpṛī} again
\end{itemize}

\begin{tcolorbox}[breakable,colback=teal!4,colframe=teal!35!black,title={Before you move on}]
What does \ur{جھونپڑی} mean? (\textquotedblleft{}A hut\textquotedblright{}.) Say it once more.
\end{tcolorbox}
\section[\texorpdfstring{\ur{احاطہ} (ahātā)}{احاطہ (ahātā)}]{\ur{احاطہ} (ahātā) — a compound; an enclosure}

\label{lesson:UR-C413-ahata}



\begin{quote}
Before the new one: say the Urdu for a large house, then the Urdu for a hut.
\end{quote}

\subsection*{You'll want to know: \ur{احاطہ}}
\textbf{\ur{احاطہ}} — \emph{ahātā} — \textquotedblleft{}a compound; an enclosure\textquotedblright{}.

\begin{quote}
\ur{کوٹھی} \ur{کے} \ur{احاطے} \ur{میں} \ur{دو} \ur{گاڑیاں} \ur{ہیں۔}
\end{quote}

\begin{quote}
\emph{koṭhī ke ahāte meṅ do gāṛiyāṅ haiṅ.}
\end{quote}

\begin{quote}
There are two cars in the compound of the house.
\end{quote}

Masculine.

\begin{cousinweb}
\textbf{\ur{احاطہ}} (\emph{ahātā}) is usually traced, through Persian, to Arabic \emph{iḥāṭa}, \textquotedblleft{}a surrounding, an enclosing\textquotedblright{}.
\end{cousinweb}

\begin{grammarlens}[title={What you've built}]
One more word for getting around.
\end{grammarlens}

\subsection*{Your turn}
\begin{itemize}
\raggedright
  \item \emph{Say it:} \emph{ahātā}
  \item \emph{Say it:} \emph{ahātā}, once more
  \item \emph{Say it:} \emph{koṭhī ke ahāte meṅ do gāṛiyāṅ haiṅ.} — There are two cars in the compound of the house.
  \item \emph{Recall:} say the Urdu for a large house, then the Urdu for a hut, then say \emph{ahātā} again
\end{itemize}

\begin{tcolorbox}[breakable,colback=teal!4,colframe=teal!35!black,title={Before you move on}]
What does \ur{احاطہ} mean? (\textquotedblleft{}A compound; an enclosure\textquotedblright{}.) Say it once more.
\end{tcolorbox}
\section[\texorpdfstring{\ur{چھپر} (chhappar)}{چھپر (chhappar)}]{\ur{چھپر} (chhappar) — a thatched roof}

\label{lesson:UR-C413-chhappar}



\begin{quote}
Before the new one: say the Urdu for a hut, then the Urdu for a compound.
\end{quote}

\subsection*{You'll want to know: \ur{چھپر}}
\textbf{\ur{چھپر}} — \emph{chhappar} — \textquotedblleft{}a thatched roof\textquotedblright{}.

\begin{quote}
\ur{جھونپڑی} \ur{کا} \ur{چھپر} \ur{بارش} \ur{میں} \ur{ٹپکتا} \ur{ہے۔}
\end{quote}

\begin{quote}
\emph{jhoṅpṛī kā chhappar bārish meṅ ṭapaktā hai.}
\end{quote}

\begin{quote}
The hut's thatch leaks in the rain.
\end{quote}

Masculine.

\begin{cousinweb}
\textbf{\ur{چھپر}} (\emph{chhappar}) is inherited: CDIAL lists the Hindi cousin \emph{chappar}, \textquotedblleft{}a sloping thatch, a hut\textquotedblright{}, under Sanskrit \emph{chattvara} (CDIAL 4976).
\end{cousinweb}

\begin{grammarlens}[title={What you've built}]
One more word for getting around.
\end{grammarlens}

\subsection*{Your turn}
\begin{itemize}
\raggedright
  \item \emph{Say it:} \emph{chhappar}
  \item \emph{Say it:} \emph{chhappar}, once more
  \item \emph{Say it:} \emph{jhoṅpṛī kā chhappar bārish meṅ ṭapaktā hai.} — The hut's thatch leaks in the rain.
  \item \emph{Recall:} say the Urdu for a hut, then the Urdu for a compound, then say \emph{chhappar} again
\end{itemize}

\begin{tcolorbox}[breakable,colback=teal!4,colframe=teal!35!black,title={Before you move on}]
What does \ur{چھپر} mean? (\textquotedblleft{}A thatched roof\textquotedblright{}.) Say it once more.
\end{tcolorbox}
\section[\texorpdfstring{\ur{بسانا} (basānā)}{بسانا (basānā)}]{\ur{بسانا} (basānā) — to found (a town); to settle (people)}

\label{lesson:UR-C413-basana}



\begin{quote}
Before the new one: say the Urdu for a compound, then the Urdu for a thatched roof.
\end{quote}

\subsection*{You'll want to know: \ur{بسانا}}
\textbf{\ur{بسانا}} — \emph{basānā} — \textquotedblleft{}to found (a town); to settle (people)\textquotedblright{}.

\begin{quote}
\ur{کسانوں} \ur{نے} \ur{دریا} \ur{کے} \ur{پاس} \ur{ایک} \ur{بستی} \ur{بسائی۔}
\end{quote}

\begin{quote}
\emph{kisānoṅ ne daryā ke pās ek bastī basāī.}
\end{quote}

\begin{quote}
Farmers founded a settlement near the river.
\end{quote}

Past, with \ur{نے} (he or she …): \ur{اس} \ur{نے} \ur{نیا} \ur{شہر} \ur{بسایا}, \emph{us ne nayā shahr basāyā}.

\begin{cousinweb}
\textbf{\ur{بسانا}} (\emph{basānā}) is built on \emph{basnā}, \textquotedblleft{}to dwell\textquotedblright{}, which CDIAL lists under Sanskrit \emph{vásati}, \textquotedblleft{}dwells\textquotedblright{} (CDIAL 11435): it is the form for making people settle; \textbf{\ur{بستی}}, \textquotedblleft{}a settlement\textquotedblright{}, belongs to the same family.
\end{cousinweb}

\begin{grammarlens}[title={What you've built}]
One more word for getting around.
\end{grammarlens}

\subsection*{Your turn}
\begin{itemize}
\raggedright
  \item \emph{Say it:} \emph{basānā}
  \item \emph{Say it:} \emph{basānā}, once more
  \item \emph{Say it:} \emph{kisānoṅ ne daryā ke pās ek bastī basāī.} — Farmers founded a settlement near the river.
  \item \emph{Recall:} say the Urdu for a compound, then the Urdu for a thatched roof, then say \emph{basānā} again
\end{itemize}

\begin{tcolorbox}[breakable,colback=teal!4,colframe=teal!35!black,title={Before you move on}]
What does \ur{بسانا} mean? (\textquotedblleft{}To found (a town); to settle (people)\textquotedblright{}.) Say it once more.
\end{tcolorbox}
